\documentclass[journal]{IEEEtran}

\usepackage[utf8]{inputenc}
\usepackage[T1]{fontenc}
\usepackage[french]{babel}
\usepackage{amsmath,amssymb,amsfonts,amsthm}
\usepackage{booktabs,multirow,array}
\usepackage{graphicx}
\usepackage{cite}
\usepackage{microtype}
\usepackage{algorithm}
\usepackage{algorithmic}
\usepackage{pgfplots}
\usepackage{tikz}
\usepackage{url}
\usepackage[colorlinks=true,linkcolor=black,citecolor=black,urlcolor=blue]{hyperref}

\usetikzlibrary{arrows.meta,positioning,shapes.geometric}
\pgfplotsset{compat=1.18}

\newtheorem{definition}{Définition}
\newtheorem{property}{Propriété}

\begin{document}

\title{Contrôle Adaptatif de Trafic Urbain par Double Deep Q-Network~: Évaluation Multi-Graines et Diagnostic Microscopique Face aux Heuristiques de Transport sous Eclipse SUMO}

\author{Mamadou~Ahmadou~Ka%
\thanks{M. A. Ka est affilié au Département Informatique, UFR Sciences et Technologies, Université Iba Der Thiam de Thiès (UIDT), Thiès, Sénégal (courriel~: \texttt{mamadou.ka@uidt.sn}).}%
}

\markboth{IEEE Transactions on Intelligent Transportation Systems,~Vol.~XX, No.~X,~Septembre~2026}%
{Ka : Contrôle Adaptatif de Trafic Urbain par Double DQN sous SUMO}

\maketitle

\begin{abstract}
La congestion récurrente aux carrefours routiers des agglomérations en développement rapide constitue un frein socio-économique majeur, exacerbé par la prédominance de feux à temporisation statique. Cet article propose un benchmark expérimental rigoureux évaluant un agent décisionnel fondé sur l'apprentissage par renforcement profond (\emph{Double Deep Q-Network}, DDQN) face aux trois paradigmes dominants du génie du trafic~: le plan fixe optimisé, la régulation actionnée par détection de lacune (\emph{actuated gap-out}) et la règle heuristique théorique Max-Pressure (Varaiya, 2013). Les expérimentations sont conduites sous le simulateur microscopique Eclipse SUMO sur la géométrie réelle d'une intersection majeure de la ville de Thiès (Sénégal), modélisée à partir de relevés OpenStreetMap. L'évaluation repose sur cinq graines stochastiques indépendantes évaluées chacune sur un horizon continu d'une heure ($3600\,\mathrm{s}$). Les résultats démontrent la supériorité de l'agent Double DQN, qui réduit le temps perdu moyen (\emph{timeLoss}) de $50{,}7\,\%$ ($15{,}5 \pm 0{,}6\,\mathrm{s}$ contre $31{,}5 \pm 0{,}5\,\mathrm{s}$ pour le plan fixe) et le temps d'attente à l'arrêt de $62{,}2\,\%$ ($9{,}0 \pm 0{,}5\,\mathrm{s}$ contre $23{,}8 \pm 0{,}4\,\mathrm{s}$), tout en maintenant le débit d'évacuation maximal ($889 \pm 21$ véhicules servis). En outre, une analyse microscopique approfondie met en lumière une pathologie critique de l'algorithme Max-Pressure naïf~: loin de souffrir d'oscillations excessives (\emph{thrashing}), celui-ci subit un verrouillage monopolistique sur l'axe majeur ($80{,}8\,\%$ du temps alloué à la Phase~2 et $0\,\%$ à la Phase~3), entraînant une famine sévère des voies secondaires et un refoulement en amont (\emph{spillback}) qui réduit le débit global de $34\,\%$. L'agent Double DQN résout spontanément cette famine en allouant $18{,}3\,\%$ du temps à la purge des mouvements secondaires, attestant d'une capacité d'arbitrage global supérieure aux règles gloutonnes locales.
\end{abstract}

\begin{IEEEkeywords}
Systèmes de Transport Intelligents (ITS), Régulation Adaptative des Feux (ATSC), Apprentissage par Renforcement Profond, Double DQN, Max-Pressure, Phénomène de Famine, Eclipse SUMO.
\end{IEEEkeywords}

\IEEEpeerreviewmaketitle

\section{Introduction}
\label{sec:intro}

\IEEEPARstart{D}{ans} les métropoles des pays émergents et particulièrement en Afrique subsaharienne, l'urbanisation accélérée s'accompagne d'une densification critique des flux motorisés sur un réseau viaire souvent contraint et peu maillé~\cite{haydari2022deep}. À Thiès (Sénégal), deuxième pôle économique du pays, les carrefours à niveau concentrent l'essentiel des congestions urbaines, provoquant retards massifs, surconsommation d'hydrocarbures et pollution atmosphérique accrue.

Historiquement, la régulation des intersections repose sur des plans de feux à temps fixe (\emph{Fixed-Time}), calibrés selon des méthodes conventionnelles telles que celle de Webster~\cite{webster1958traffic}. Bien que robustes et simples d'exploitation, ces plans statiques sont structurellement incapables de s'ajuster aux fluctuations stochastiques de la demande (pics asymétriques, incidents de voirie, arrivées en peloton). Les systèmes industriels adaptatifs coordonnés, tels que SCATS~\cite{lowrie1990scats} ou SCOOT~\cite{hunt1981scoot}, requièrent quant à eux une instrumentation lourde et intrusive de la chaussée (boucles magnétiques inductives tous les 50 mètres, réseaux de transmission filaires dédiés) ainsi que des coûts récurrents de maintenance prohibitifs pour les collectivités locales.

Ce constat a suscité un engouement croissant pour les approches fondées sur l'apprentissage par renforcement profond (\emph{Deep Reinforcement Learning}, DRL)~\cite{wei2018intellilight,wei2019presslight,chen2020toward,zhang2020pdlight}. En observant l'état local du trafic via des capteurs visuels peu onéreux (caméras intelligentes \emph{Edge AI}), un agent autonome peut apprendre une politique optimale d'allocation des feux sans nécessiter de modèle paramétrique explicite de la demande.

Cependant, comme le soulignent plusieurs revues critiques récentes de la littérature~\cite{wei2021survey,zhao2024survey}, une majorité écrasante des contributions récentes en DRL pèche par un manque de rigueur méthodologique~:
\begin{enumerate}
    \item \textbf{Baselines limitées~:} Les performances de l'agent ne sont que trop souvent comparées au seul plan fixe, ignorant les contrôleurs actionnés (\emph{actuated control}) et les heuristiques théoriques dominantes de la recherche en transport, en premier lieu l'algorithme Max-Pressure formulé par Varaiya~\cite{varaiya2013max}.
    \item \textbf{Absence de diagnostic comportemental microscopique~:} L'analyse se borne fréquemment à rapporter une courbe de récompense cumulative ou un gain macroscopique global, sans disséquer les mécanismes sous-jacents de sélection de phase (durée de vert effectif, fréquence de commutation, gestion des tournes critiques).
    \item \textbf{Évaluation sur environnements idéalisés~:} Les topologies de test sont souvent réduites à des grilles 2D symétriques et synthétiques, masquant les effets de goulot d'étranglement propres aux carrefours réels asymétriques.
\end{enumerate}

Pour combler ces lacunes, cet article propose une étude comparative et comportementale rigoureuse. Nos contributions se déclinent comme suit~:
\begin{itemize}
    \item Nous formulons le contrôle d'une intersection réelle à 8 voies de la ville de Thiès (Sénégal) sous la forme d'un Processus de Décision Markovien (MDP) compact à 13 dimensions, avec une fonction de récompense différentielle sans pénalité arbitraire.
    \item Nous déployons et entraînons un agent \textbf{Double DQN} (DDQN) intégrant des masques de contraintes sécuritaires (verts minimum et maximum, phases jaunes d'inter-vert).
    \item Nous confrontons cet agent à trois familles de contrôleurs~: le plan fixe de référence (Webster), le régulateur actionné par détection de trou de trafic (\emph{gap-out}) et deux variantes d'agrégation de l'algorithme Max-Pressure (\emph{Sum} et \emph{Max}).
    \item Nous démontrons, par une analyse microscopique des trajectoires de simulation sous Eclipse SUMO, que le dysfonctionnement de Max-Pressure n'est pas imputable à des oscillations intempestives (\emph{thrashing}), mais à un phénomène pernicieux de \textbf{famine de phase} (\emph{phase starvation}) conduisant à un refoulement amont (\emph{spillback}) que l'agent Double DQN résout spontanément.
\end{itemize}

Le reste de cet article est structuré comme suit. La Section~\ref{sec:related} passe en revue l'état de l'art. La Section~\ref{sec:formulation} expose la formulation mathématique du problème et des contrôleurs. La Section~\ref{sec:algo} détaille l'algorithme Double DQN et son architecture neuronale. La Section~\ref{sec:setup} décrit le protocole expérimental sous SUMO. La Section~\ref{sec:results} présente les résultats quantitatifs multi-graines et l'analyse microscopique de la famine. Enfin, la Section~\ref{sec:conclusion} conclut et ouvre sur les perspectives.

\section{État de l'Art}
\label{sec:related}

\subsection{Méthodes Traditionnelles et Heuristiques de Transport}
La méthode de Webster~\cite{webster1958traffic} constitue le socle du dimensionnement des cycles de feux à temps fixe. Elle calcule la durée optimale de cycle $C^*$ en minimisant le retard total sous l'hypothèse d'arrivées poissoniennes stationnaires. En présence de variabilité intrajournalière, la régulation actionnée (\emph{Actuated Control}) ajuste dynamiquement le maintien du feu vert par détection locale~: tant que les véhicules se succèdent à un intervalle inférieur à un seuil de passage critique ($\Delta t_{\text{gap}}$), le vert est prolongé jusqu'à concurrence d'un plafond $g_{\max}$~\cite{day2016detector}.

Sur le plan théorique, l'algorithme \textbf{Max-Pressure} (MP) introduit par Varaiya~\cite{varaiya2013max} et étendu par Lioris \emph{et al.}~\cite{lioris2016max} a marqué un tournant. Inspiré du contrôle par contre-pression dans les réseaux de télécommunications, MP prouve mathématiquement qu'en sélectionnant à chaque instant la phase qui maximise la différence de files d'attente entre les voies amont et aval, le réseau de trafic demeure globalement stable (les files n'explosent pas à l'infini) dès lors que la demande se situe dans la région de capacité admissible (\emph{capacity region}). Cependant, comme nous le mettrons en évidence dans cette étude, les garanties asymptotiques de Max-Pressure reposent sur des hypothèses continues (\emph{store-and-forward}) qui se révèlent vulnérables en environnement microscopique discret à carrefour isolé lorsque les longueurs de stockage et les structures de voies sont asymétriques.

\subsection{Apprentissage par Renforcement Profond pour l'ATSC}
L'avènement du Deep Q-Network (DQN)~\cite{mnih2015human} a catalysé son application au contrôle adaptatif des feux (\emph{Adaptive Traffic Signal Control}, ATSC)~\cite{genders2016using}. Wei \emph{et al.} ont posé les fondations du domaine avec \textbf{IntelliLight}~\cite{wei2018intellilight}, démontrant qu'un agent DQN alimenté par des états continus (files, vitesses, positions) surpasse significativement les heuristiques traditionnelles sur des intersections réelles.

Pour combler le fossé entre la théorie du transport et le DRL, Wei \emph{et al.} ont proposé \textbf{PressLight}~\cite{wei2019presslight}, intégrant la formulation de la pression de Varaiya au sein même de la fonction de récompense de l'agent. Parallèlement, Chen \emph{et al.} ont développé \textbf{MPLight}~\cite{chen2020toward}, combinant Max-Pressure et factorisation d'état pour le contrôle décentralisé à grande échelle. Récemment, Zhang \emph{et al.} ont introduit \textbf{PDLight}~\cite{zhang2020pdlight}, modifiant la mesure de pression pour intégrer explicitement la capacité résiduelle des voies d'évacuation aval (\emph{Pressure with Remaining Capacity of Outgoing Lane}, PRCOL) afin de juguler les blocages par refoulement. 

Notre étude se distingue de ces travaux par une analyse comportementale rigoureuse d'un agent Double DQN face aux variantes de Max-Pressure sur un carrefour africain réel sous SUMO, en décortiquant les micro-mécanismes de famine temporelle et spatiale.

\section{Formulation du Problème et des Modèles}
\label{sec:formulation}

\subsection{Topologie de l'Intersection et Notations}
Considérons une intersection routière isolée à niveau (Figure~\ref{fig:topology_graph}). Soit $\mathcal{L}_{\text{in}} = \{l_1, \dots, l_M\}$ l'ensemble des voies d'entrée amont et $\mathcal{L}_{\text{out}} = \{l'_1, \dots, l'_N\}$ l'ensemble des voies de sortie aval. Un mouvement de trafic $m = (l, l')$ représente la trajectoire autorisée d'une voie entrante $l \in \mathcal{L}_{\text{in}}$ vers une voie sortante $l' \in \mathcal{L}_{\text{out}}$.

\begin{figure}[t]
\centering
\begin{tikzpicture}[node distance=1.4cm, auto, >=Latex, thick]
    \node [draw, rectangle, fill=blue!10, rounded corners, minimum width=2.8cm, minimum height=0.9cm, align=center] (in_north) {\small \textbf{Approche Nord} \\ \footnotesize 2 voies (Ligne 1)};
    \node [draw, rectangle, fill=blue!10, rounded corners, minimum width=2.8cm, minimum height=0.9cm, align=center, below=1.8cm of in_north] (in_south) {\small \textbf{Approche Sud} \\ \footnotesize 2 voies (Ligne 2)};
    \node [draw, rectangle, fill=blue!10, rounded corners, minimum width=2.2cm, minimum height=0.9cm, align=center, left=1.2cm of in_north, yshift=-1.35cm] (in_west) {\small \textbf{Approche Ouest} \\ \footnotesize 2 voies};
    \node [draw, rectangle, fill=blue!10, rounded corners, minimum width=2.2cm, minimum height=0.9cm, align=center, right=1.2cm of in_north, yshift=-1.35cm] (in_east) {\small \textbf{Approche Est} \\ \footnotesize 2 voies};

    \node [draw, circle, fill=red!20, inner sep=2pt, minimum size=1.3cm, align=center] (junction) at ([yshift=-0.9cm]in_north.south) {\textbf{Junction} \\ \scriptsize $J_0$};

    \draw [<->, color=black!70] (in_north) -- (junction);
    \draw [<->, color=black!70] (in_south) -- (junction);
    \draw [<->, color=black!70] (in_west) -- (junction);
    \draw [<->, color=black!70] (in_east) -- (junction);
\end{tikzpicture}
\caption{Topologie schématique de l'intersection de Thiès (8 voies d'entrée réparties sur 4 branches cardinales avec mouvements directs et tourne-à-gauche protégés).}
\label{fig:topology_graph}
\end{figure}

L'intersection est régulée par un ensemble fini de phases de feux verts non conflictuelles $\mathcal{P} = \{p_0, p_1, p_2, p_3\}$, où chaque phase $p \in \mathcal{P}$ autorise simultanément un sous-ensemble de mouvements non sécants $\mathcal{M}(p)$~:
\begin{itemize}
    \item \textbf{Phase 0~:} Circulation directe sur l'axe majeur (Nord-Sud) ;
    \item \textbf{Phase 1~:} Tourne-à-gauche protégé sur l'axe majeur ;
    \item \textbf{Phase 2~:} Circulation directe sur l'axe secondaire (Est-Ouest) ;
    \item \textbf{Phase 3~:} Tourne-à-gauche protégé sur l'axe secondaire.
\end{itemize}
Toute transition d'une phase active $p$ vers une phase distincte $p'$ impose l'insertion systématique d'une phase de dégagement jaune de durée $y = 3\,\mathrm{s}$.

\subsection{Formulation de l'Heuristique Max-Pressure}
Dans le cadre d'un carrefour isolé où l'évacuation sur les voies aval $\mathcal{L}_{\text{out}}$ n'est pas entravée par une intersection adjacente, la pression $P_m(t)$ d'un mouvement $m$ correspond directement au nombre de véhicules arrêtés ou en attente $x_l(t)$ sur sa voie entrante d'origine $l$~:
\begin{equation}
P_m(t) = x_l(t), \quad \forall m = (l, l').
\end{equation}
La pression globale d'une phase candidate $p \in \mathcal{P}$ est alors calculée par agrégation des pressions de l'ensemble de ses mouvements autorisés $\mathcal{M}(p)$. Nous évaluons deux fonctions d'agrégation~:
\begin{equation}
\label{eq:mp_sum}
\text{Mode Somme (Sum) : } P_{\text{sum}}(p, t) = \sum_{m \in \mathcal{M}(p)} x_{l(m)}(t),
\end{equation}
\begin{equation}
\label{eq:mp_max}
\text{Mode Maximum (Max) : } P_{\text{max}}(p, t) = \max_{m \in \mathcal{M}(p)} x_{l(m)}(t).
\end{equation}
À chaque instant de décision $t$, l'algorithme Max-Pressure sélectionne la phase d'indice~:
\begin{equation}
a_{\text{MP}}^*(t) = \arg\max_{p \in \mathcal{A}_{\text{valid}}(t)} P(p, t),
\end{equation}
où $\mathcal{A}_{\text{valid}}(t)$ restreint le choix de l'action afin de garantir le respect strict du temps de vert minimal de sécurité ($g_{\min} = 10\,\mathrm{s}$) et du plafond maximal ($g_{\max} = 60\,\mathrm{s}$).

\subsection{Formulation en Processus de Décision Markovien (MDP)}
Nous modélisons le problème de contrôle sous la forme d'un Processus de Décision Markovien défini par le quintuplet $\langle \mathcal{S}, \mathcal{A}, \mathcal{P}_{ss'}, \mathcal{R}, \gamma \rangle$~:

\subsubsection{Espace des États (Observation)}
Le vecteur d'observation $\mathbf{s}_t \in \mathbb{R}^{13}$ capturé par l'agent à l'instant de décision $t$ est structuré en quatre sous-groupes d'attributs continus normalisés dans l'intervalle $[0, 1]$~:
\begin{equation}
\mathbf{s}_t = \big[ \mathbf{q}_t^{\top}, \, \mathbf{k}_t^{\top}, \, \mathbf{e}_t^{\top}, \, \delta_t \big]^{\top} \in \mathbb{R}^{13},
\end{equation}
où~:
\begin{itemize}
    \item $\mathbf{q}_t = [q_0, q_1, q_2, q_3]^{\top} \in [0, 1]^4$ quantifie les longueurs de file normalisées sur les quatre approches d'entrée, où $q_i = \min(n_{\text{stop}, i} / Q_{\max}, 1{,}0)$ avec un seuil de saturation fixé à $Q_{\max} = 20$ véhicules arrêtés ($v < 0{,}1\,\mathrm{m/s}$).
    \item $\mathbf{k}_t = [k_0, k_1, k_2, k_3]^{\top} \in [0, 1]^4$ mesure les densités spatiales normalisées d'occupation de la chaussée, $k_i = n_{\text{tot}, i} / K_{\text{cap}, i}$, où $K_{\text{cap}, i} = L_i / 7{,}5\,\mathrm{m}$ représente la capacité physique théorique de stockage de la voie $i$ de longueur $L_i$.
    \item $\mathbf{e}_t \in \{0, 1\}^4$ est un encodage \emph{One-Hot} désignant la phase verte actuellement active ($e_j = 1$ si la phase $j$ est au vert, $0$ sinon).
    \item $\delta_t \in \{0, 1\}$ est un booléen de garde sécuritaire valant $1$ si la phase en cours a atteint la durée minimale réglementaire ($t_{\text{green}} \ge g_{\min}$), autorisant une transition, et $0$ dans le cas contraire.
\end{itemize}

\subsubsection{Espace des Actions}
L'espace des actions discrètes est défini par l'indexation directe des phases vertes autorisées~:
\begin{equation}
\mathcal{A} = \{0, 1, 2, 3\},
\end{equation}
où choisir $a_t = p$ ordonne à l'automate d'engager ou de maintenir la phase $p$. Si $a_t \neq a_{t-1}$, le simulateur déclenche automatiquement une phase de transition jaune de $y = 3\,\mathrm{s}$ avant d'activer le vert de la phase $a_t$.

\subsubsection{Fonction de Récompense}
La fonction de récompense est conçue pour pénaliser directement l'accumulation de retards sans introduire de pondérations heuristiques instables. Elle est définie comme la variation différentielle du temps d'attente cumulé global $W(t) = \sum_{v \in \mathcal{V}(t)} w_v(t)$ de l'ensemble des véhicules présents sur l'intersection~:
\begin{equation}
\label{eq:reward}
r_t = \frac{W(t) - W(t + \Delta t_{\text{step}})}{100}.
\end{equation}
Cette formulation confère deux propriétés mathématiques fondamentales~:
\begin{property}
Toute résorption nette du temps d'attente ($W(t + \Delta t) < W(t)$) produit une récompense positive, alignant l'agent sur la maximisation du débit de sortie.
\end{property}
\begin{property}
Le coût d'inter-vert induit par les phases jaunes se trouve naturellement pénalisé dans la fonction de valeur sans artifice externe~: commuter de phase bloque temporairement le carrefour pendant $3\,\mathrm{s}$, induisant une hausse transitoire de $W(t)$ qui pénalise mécaniquement les changements de phase futiles.
\end{property}

\section{Algorithme Double DQN et Architecture}
\label{sec:algo}

\subsection{Découplage de Bellman et Réseau de Valeur}
Dans le cadre classique du Q-learning profond (DQN)~\cite{mnih2015human}, le biais de maximisation systématique $\max_{a'} Q(s', a'; \boldsymbol{\theta})$ tend à surestimer sévèrement les valeurs $Q$, entraînant une dégradation des politiques en environnement stochastique. L'algorithme Double DQN~\cite{van2016deep} élimine ce biais en découplant la sélection de l'action optimale de son évaluation~:
\begin{equation}
\label{eq:target_double_dqn}
Y_t^{\text{DoubleQ}} = r_t + \gamma \, Q\big(s_{t+1}, \, \arg\max_{a'} Q(s_{t+1}, a'; \boldsymbol{\theta}_t); \, \boldsymbol{\theta}_t^-\big),
\end{equation}
où $\boldsymbol{\theta}_t$ désigne les paramètres du réseau de neurones en ligne (\emph{Online Network}) et $\boldsymbol{\theta}_t^-$ représente les poids du réseau cible (\emph{Target Network}), gelés et mis à jour périodiquement toutes les $K_{\text{target}}$ itérations.

L'optimisation des paramètres $\boldsymbol{\theta}$ est réalisée par descente de gradient stochastique sur la fonction de perte quadratique moyenne (MSE)~:
\begin{equation}
\label{eq:loss}
\mathcal{L}(\boldsymbol{\theta}) = \mathbb{E}_{(s, a, r, s') \sim \mathcal{D}} \left[ \left( Y^{\text{DoubleQ}} - Q(s, a; \boldsymbol{\theta}) \right)^2 \right],
\end{equation}
où $\mathcal{D}$ désigne le tampon d'expérience (\emph{Replay Buffer}) d'une capacité de $N_{\text{buffer}} = 50\,000$ transitions.

\subsection{Architecture Neuronale et Masquage d'Actions}
L'architecture du réseau de valeur est un Perceptron Multicouche (MLP) fully-connected compact et efficient, adapté au calcul embarqué sur microprocesseur de voirie~:
\begin{align}
\mathbf{h}_1 &= \text{ReLU}(\mathbf{W}_1 \mathbf{s} + \mathbf{b}_1), \quad \mathbf{W}_1 \in \mathbb{R}^{64 \times 13}, \\
\mathbf{h}_2 &= \text{ReLU}(\mathbf{W}_2 \mathbf{h}_1 + \mathbf{b}_2), \quad \mathbf{W}_2 \in \mathbb{R}^{64 \times 64}, \\
\mathbf{Q}(\mathbf{s}) &= \mathbf{W}_3 \mathbf{h}_2 + \mathbf{b}_3, \quad \mathbf{W}_3 \in \mathbb{R}^{4 \times 64}.
\end{align}

Pour garantir l'intégrité physique de la circulation et interdire tout comportement erratique, un masque d'admissibilité des actions $\mathcal{M}(s_t) \subset \mathcal{A}$ est appliqué avant l'étape $\varepsilon$-greedy~:
\begin{equation}
\mathcal{M}(s_t) = \begin{cases}
\{a_{t-1}\} & \text{si } t_{\text{green}} < g_{\min}, \\
\mathcal{A} \setminus \{a_{t-1}\} & \text{si } t_{\text{green}} \ge g_{\max}, \\
\mathcal{A} & \text{sinon.}
\end{cases}
\end{equation}
L'algorithme d'apprentissage et de contrôle est synthétisé dans l'Algorithme~\ref{alg:double_dqn}.

\begin{algorithm}[t]
\caption{Contrôle Adaptatif par Double DQN sous Contraintes}
\label{alg:double_dqn}
\begin{algorithmic}[1]
\REQUIRE Réseau en ligne $Q_{\boldsymbol{\theta}}$, réseau cible $Q_{\boldsymbol{\theta}^-}$, Replay Buffer $\mathcal{D}$.
\REQUIRE Hyperparamètres~: taux d'apprentissage $\alpha = 10^{-3}$, facteur d'actualisation $\gamma = 0{,}95$, pas de décision $\Delta t = 5\,\mathrm{s}$.
\STATE Initialiser $\boldsymbol{\theta}^- \leftarrow \boldsymbol{\theta}$, $\mathcal{D} \leftarrow \emptyset$.
\FOR{chaque épisode $= 1$ \TO $E$}
    \STATE Réinitialiser l'environnement SUMO, observer $\mathbf{s}_0 \in \mathbb{R}^{13}$.
    \FOR{$t = 0, \Delta t, 2\Delta t, \dots, T$}
        \STATE Déterminer le masque d'actions valides $\mathcal{M}(\mathbf{s}_t)$.
        \STATE Avec probabilité $\varepsilon$, sélectionner $a_t \in \mathcal{M}(\mathbf{s}_t)$ aléatoirement ;
        \STATE Sinon, sélectionner $a_t = \arg\max_{a \in \mathcal{M}(\mathbf{s}_t)} Q(\mathbf{s}_t, a; \boldsymbol{\theta})$.
        \IF{$a_t \neq a_{t-1}$}
            \STATE Exécuter phase jaune $y = 3\,\mathrm{s}$ via TraCI.
        \ENDIF
        \STATE Maintenir la phase $a_t$ pendant $\Delta t$, simuler le trafic.
        \STATE Observer le nouvel état $\mathbf{s}_{t+1}$ et calculer $r_t$ via~(\ref{eq:reward}).
        \STATE Stocker la transition $(\mathbf{s}_t, a_t, r_t, \mathbf{s}_{t+1})$ dans $\mathcal{D}$.
        \STATE Échantillonner un mini-lot $B \sim \mathcal{D}$ de taille 64.
        \STATE Calculer les cibles $Y^{\text{DoubleQ}}$ via~(\ref{eq:target_double_dqn}).
        \STATE Mettre à jour $\boldsymbol{\theta}$ par descente de gradient sur $\mathcal{L}(\boldsymbol{\theta})$~(\ref{eq:loss}).
        \STATE Périodiquement, synchroniser $\boldsymbol{\theta}^- \leftarrow \boldsymbol{\theta}$.
    \ENDFOR
\ENDFOR
\end{algorithmic}
\end{algorithm}

\begin{figure}[t]
\centering
\begin{tikzpicture}
\begin{axis}[
    width=\columnwidth,
    height=4.6cm,
    xlabel={Épisode d'apprentissage},
    ylabel={Longueur de file moyenne (véh.)},
    xmin=1, xmax=150,
    ymin=1.5, ymax=11.0,
    grid=major,
    grid style={dashed,gray!25},
    axis lines=left,
    legend pos=north east,
    legend style={font=\footnotesize, fill=white, fill opacity=0.8}
]
\addplot[thick, color=blue!80, mark=*, mark size=1.1pt] coordinates {
    (1,8.45) (6,9.77) (11,9.88) (16,8.81) (21,7.09) (26,6.73) (31,6.17) (36,4.82) (41,4.19) (46,4.01) (51,3.80) (56,3.40) (61,3.22) (66,3.14) (71,3.08) (76,3.03) (81,3.02) (86,2.94) (91,2.82) (96,2.77) (101,2.76) (106,2.72) (111,2.68) (116,2.65) (121,2.58) (126,2.62) (131,2.55) (136,2.48) (141,2.54) (146,2.50)
};
\legend{Moyenne mobile (10 ép.)}
\end{axis}
\end{tikzpicture}
\caption{Convergence de la politique Double DQN sur 150 épisodes d'entraînement sous SUMO (réduction de la file moyenne de 9,88 à 2,50 véhicules).}
\label{fig:training_curve}
\end{figure}

\section{Protocole Expérimental sous Eclipse SUMO}
\label{sec:setup}

\subsection{Environnement de Simulation et Réseau Routier}
L'ensemble des expérimentations est exécuté sous la plateforme microscopique open-source \textbf{Eclipse SUMO} (version 1.26.0) interfacée en temps réel via l'API socket Python \textbf{TraCI} (\emph{Traffic Control Interface}). 

La topologie modélisée reproduit fidèlement le carrefour principal de la ville de \textbf{Thiès (Sénégal)}, extrait depuis la base cartographique OpenStreetMap et compilé avec l'utilitaire \texttt{netconvert}. L'intersection comprend quatre branches d'accès composées chacune de 2 voies de circulation (largeur $3{,}25\,\mathrm{m}$, vitesse limite $50\,\mathrm{km/h}$), soit 8 voies d'entrée et 8 voies de sortie convergentes sur le nœud central $J_0$.

Le pas de calcul physique de la simulation est paramétré à $\Delta t_{\text{sim}} = 0{,}1\,\mathrm{s}$ avec le modèle de poursuite de Krauß. Pour éliminer tout artefact induit par les automates internes de SUMO, le contrôle des feux est opéré en mode exclusif via la commande native TraCI \texttt{setRedYellowGreenState()}. L'agent intervient à un intervalle de décision régulier fixé à $\Delta t_{\text{decision}} = 5{,}0\,\mathrm{s}$.

\subsection{Demande de Trafic et Graines d'Évaluation}
Afin d'évaluer la robustesse statistique des politiques, nous avons défini un protocole multi-graines rigoureux~:
\begin{itemize}
    \item Les épisodes d'évaluation durent $T_{\text{eval}} = 3600\,\mathrm{s}$ (soit 1 heure d'observation continue du trafic urbain).
    \item L'évaluation est conduite sur cinq graines stochastiques indépendantes~:
    \begin{equation}
    \mathcal{S}_{\text{eval}} = \{7, 13, 21, 42, 77\}.
    \end{equation}
    Ces graines régissent les tirages pseudo-aléatoires des instants d'injection et des itinéraires des véhicules dans le fichier de configuration \texttt{rl.sumocfg}. Elles sont strictement disjointes des graines utilisées lors de la phase d'apprentissage du modèle DDQN.
    \item À l'évaluation, la politique DDQN est exploitée de façon purement déterministe ($\varepsilon = 0$), sans exploration.
\end{itemize}

\subsection{Contrôleurs de Référence (Baselines)}
Quatre contrôleurs de référence sont déployés dans des conditions expérimentales rigoureusement identiques~:
\begin{enumerate}
    \item \textbf{Plan Fixe (Webster) : } Programme cyclique à 4 phases pré-dimensionné selon les flux moyens observés~: Phase~0 ($27\,\mathrm{s}$), Jaune ($3\,\mathrm{s}$), Phase~1 ($6\,\mathrm{s}$), Jaune ($3\,\mathrm{s}$), Phase~2 ($27\,\mathrm{s}$), Jaune ($3\,\mathrm{s}$), Phase~3 ($6\,\mathrm{s}$), Jaune ($3\,\mathrm{s}$), soit une durée totale de cycle de $C = 78\,\mathrm{s}$.
    \item \textbf{Contrôleur Actionné (\emph{Actuated}) : } Algorithme basé sur la détection locale de lacune (\emph{gap-out}). Le vert en cours est prolongé tant que l'intervalle temporel entre véhicules successifs reste inférieur à $\Delta t_{\text{gap}} = 3{,}0\,\mathrm{s}$, sous réserve des bornes $g_{\min} = 10\,\mathrm{s}$ et $g_{\max} = 60\,\mathrm{s}$.
    \item \textbf{Max-Pressure Naïf (\emph{MP-Sum}) : } Règle heuristique de Varaiya sélectionnant la phase maximisant la somme non pondérée des files entrantes selon~(\ref{eq:mp_sum}).
    \item \textbf{Max-Pressure Normalisé (\emph{MP-Max}) : } Règle sélectionnant la phase dont la file d'approche maximale est la plus élevée selon~(\ref{eq:mp_max}).
\end{enumerate}

\subsection{Métriques d'Évaluation de Performance}
Pour chaque véhicule $v \in \mathcal{V}_{\text{finished}}$ ayant complété son trajet durant l'heure de simulation, les métriques microscopiques enregistrées dans le fichier \texttt{tripinfo.xml} sont agrégées~:
\begin{itemize}
    \item \textbf{Temps perdu moyen (\emph{timeLoss})}~: Écart entre le temps de trajet réel et le temps en vitesse libre théorique sans congestion ($s$).
    \item \textbf{Temps d'attente moyen (\emph{waitingTime})}~: Temps cumulé passé à l'arrêt complet ($v < 0{,}1\,\mathrm{m/s}$) ($s$).
    \item \textbf{Durée moyenne de trajet (\emph{duration})}~: Temps total de parcours dans le réseau ($s$).
    \item \textbf{Vitesse moyenne d'écoulement ($v$)}~: Rapport distance parcourue sur durée moyenne ($\mathrm{m/s}$).
    \item \textbf{Débit d'évacuation ($n$)}~: Nombre total de véhicules parvenus à destination durant les $3600\,\mathrm{s}$.
    \item \textbf{Véhicules bloqués en amont ($N_{\text{block}}$)}~: Véhicules générés par SUMO mais incapables de s'insérer sur le réseau par manque d'espace libre (\emph{teleport/insertion spillback}).
\end{itemize}

\begin{table*}[t]
\centering
\caption{Performances comparatives consolidées des contrôleurs sur 5 graines stochastiques indépendantes ($3600\,\mathrm{s}$ d'horizon)}
\label{tab:benchmark_results}
\vspace{0.15cm}
\begin{tabular}{lcccccc}
\toprule
\textbf{Contrôleur} & \textbf{timeLoss (s)} $\downarrow$ & \textbf{waitingTime (s)} $\downarrow$ & \textbf{Durée de trajet (s)} $\downarrow$ & \textbf{Vitesse (m/s)} $\uparrow$ & \textbf{Débit achevé ($n$)} $\uparrow$ & \textbf{Variation / Fixe} \\
\midrule
\textbf{Plan Fixe (Webster)} & $31{,}5 \pm 0{,}5$ & $23{,}8 \pm 0{,}4$ & $48{,}5 \pm 0{,}6$ & $6{,}37 \pm 0{,}07$ & $872 \pm 23$ & Référence ($0\,\%$) \\
\textbf{Actionné (\emph{Actuated})} & $34{,}2 \pm 1{,}1$ & $25{,}8 \pm 0{,}8$ & $51{,}2 \pm 1{,}2$ & $5{,}72 \pm 0{,}11$ & $874 \pm 22$ & $+8{,}7\,\%$ \\
\textbf{Max-Pressure (\emph{MP-Sum})} & $71{,}7 \pm 15{,}2$ & $57{,}3 \pm 12{,}6$ & $86{,}6 \pm 14{,}8$ & $4{,}12 \pm 0{,}45$ & $580 \pm 28$ & $+127{,}7\,\%$ \\
\textbf{Max-Pressure (\emph{MP-Max})} & $45{,}7 \pm 3{,}2$ & $37{,}7 \pm 2{,}9$ & $63{,}2 \pm 3{,}5$ & $6{,}65 \pm 0{,}18$ & $798 \pm 18$ & $+45{,}1\,\%$ \\
\midrule
\textbf{Double DQN (Notre approche)} & $\mathbf{15{,}5 \pm 0{,}6}$ & $\mathbf{9{,}0 \pm 0{,}5}$ & $\mathbf{32{,}7 \pm 0{,}7}$ & $\mathbf{8{,}16 \pm 0{,}09}$ & $\mathbf{889 \pm 21}$ & $\mathbf{-50{,}7\,\%}$ \\
\bottomrule
\end{tabular}
\end{table*}

\section{Résultats Expérimentaux et Analyses}
\label{sec:results}

\subsection{Analyse Macro-Performances Globale}
Le Tableau~\ref{tab:benchmark_results} présente la synthèse numérique des métriques de trafic sur les 5 graines stochastiques d'évaluation. La Figure~\ref{fig:bar_comparison} illustre graphiquement les disparités observées en termes de temps perdu et de temps d'attente.

\begin{figure}[b]
\centering
\begin{tikzpicture}
\begin{axis}[
    ybar,
    bar width=11pt,
    width=\columnwidth,
    height=5.8cm,
    ylabel={Temps moyen par véhicule (s)},
    symbolic x coords={Fixe, Actionné, MP-Sum, MP-Max, DDQN},
    xtick=data,
    nodes near coords,
    nodes near coords align={vertical},
    every node near coord/.append style={font=\tiny},
    ymin=0, ymax=85,
    grid=major,
    grid style={dashed,gray!25},
    axis lines=left,
    legend pos=north west,
    legend style={font=\footnotesize, fill=white, fill opacity=0.8}
]
\addplot[fill=blue!75, draw=blue!90] coordinates {
    (Fixe, 31.5) (Actionné, 34.2) (MP-Sum, 71.7) (MP-Max, 45.7) (DDQN, 15.5)
};
\addplot[fill=red!75, draw=red!90] coordinates {
    (Fixe, 23.8) (Actionné, 25.8) (MP-Sum, 57.3) (MP-Max, 37.7) (DDQN, 9.0)
};
\legend{\textit{timeLoss}, \textit{waitingTime}}
\end{axis}
\end{tikzpicture}
\caption{Comparaison du temps perdu (\emph{timeLoss}) et du temps d'attente à l'arrêt (\emph{waitingTime}) selon les différents contrôleurs testés (moyennes sur 5 graines).}
\label{fig:bar_comparison}
\end{figure}

Les enseignements majeurs suivants se dégagent des observations~:
\begin{itemize}
    \item \textbf{Efficacité remarquable du Double DQN~:} L'agent appris surpasse massivement tous les régulateurs conventionnels. Il divise le temps perdu par deux par rapport au plan fixe ($-50{,}7\,\%$, soit $15{,}5\,\mathrm{s}$ vs $31{,}5\,\mathrm{s}$) et abaisse le temps d'attente à seulement $9{,}0\,\mathrm{s}$ (gain de $-62{,}2\,\%$). La vitesse moyenne d'écoulement atteint $8{,}16\,\mathrm{m/s}$, soit un gain de fluidité de $+28{,}1\,\%$.
    \item \textbf{Stabilité et répétabilité~:} L'écart-type inter-graines sous DDQN reste extrêmement contenu ($\pm 0{,}6\,\mathrm{s}$ sur le \emph{timeLoss}), prouvant que la politique a acquis une dynamique robuste indépendante des fluctuations d'injection.
    \item \textbf{Faiblesse relative du mode Actionné~:} Le régulateur actionné enregistre une légère détérioration ($+8{,}7\,\%$ de temps perdu). Ce phénomène est classique aux carrefours denses à plusieurs voies~: en prolongeant le vert au profit de quelques véhicules isolés sur l'axe prioritaire jusqu'au plafond $g_{\max}$, il pénalise lourdement les axes transverses en attente.
    \item \textbf{Dégradation catastrophique de Max-Pressure Naïf~:} L'algorithme MP-Sum subit une explosion de son temps perdu ($71{,}7 \pm 15{,}2\,\mathrm{s}$, soit $+127{,}7\,\%$) et une chute dramatique de son débit d'évacuation ($580$ véhicules achevés contre $889$ pour le DDQN, soit $-34{,}8\,\%$ de trafic non évacué).
\end{itemize}

\subsection{Diagnostic Microscopique : Disséquer l'Échec de Max-Pressure}
\label{subsec:microscopic}

L'effondrement des performances de Max-Pressure en environnement microscopique exigeait une investigation comportementale minutieuse. Une première hypothèse intuitive aurait été d'attribuer cet échec à un phénomène d'oscillations frénétiques (\emph{thrashing}), où l'algorithme basculerait continuellement d'une phase à l'autre dès qu'un véhicule se présente, gaspillant le temps utile en phases jaunes de dégagement ($3\,\mathrm{s}$).

Pour tester rigoureusement cette hypothèse, nous avons instrumenté la simulation sous la graine 7 ($3600\,\mathrm{s}$) en consignant le nombre total de commutations de phases (\emph{phase switches}) et la répartition exacte du temps de vert cumulé alloué à chaque phase. Les résultats sont présentés dans le Tableau~\ref{tab:diagnostic_phases} et visualisés sur la Figure~\ref{fig:stacked_distribution}.

\begin{table}[t]
\centering
\caption{Diagnostic microscopique d'allocation des phases sur la graine 7 ($3600\,\mathrm{s}$)}
\label{tab:diagnostic_phases}
\vspace{0.15cm}
\resizebox{\columnwidth}{!}{%
\begin{tabular}{lcccccc}
\toprule
\textbf{Contrôleur} & \textbf{Commut.} & \textbf{Ph.~0} & \textbf{Ph.~1} & \textbf{Ph.~2} & \textbf{Ph.~3} & \textbf{Bloqués amont} \\
\midrule
\textbf{Plan Fixe} & 159 & $40{,}9\,\%$ & $9{,}1\,\%$ & $40{,}9\,\%$ & $9{,}1\,\%$ & 0 \\
\textbf{Double DQN} & 210 & $41{,}5\,\%$ & $0{,}0\,\%$ & $40{,}1\,\%$ & $\mathbf{18{,}3\,\%}$ & \textbf{0} \\
\textbf{MP-Sum} & \textbf{91} & $19{,}2\,\%$ & $0{,}0\,\%$ & $\mathbf{80{,}8\,\%}$ & $\mathbf{0{,}0\,\%}$ & \textbf{207} \\
\bottomrule
\end{tabular}%
}
\end{table}

\begin{figure}[t]
\centering
\begin{tikzpicture}
\begin{axis}[
    ybar stacked,
    bar width=16pt,
    width=0.95\columnwidth,
    height=5.2cm,
    ylabel={Répartition du vert (\%)},
    symbolic x coords={Plan Fixe, Double DQN, MP-Sum},
    xtick=data,
    ymin=0, ymax=100,
    grid=major,
    grid style={dashed,gray!25},
    axis lines=left,
    legend pos=outer north east,
    legend style={font=\tiny, fill=white, fill opacity=0.8}
]
\addplot[fill=blue!80] coordinates {(Plan Fixe, 40.9) (Double DQN, 41.5) (MP-Sum, 19.2)};
\addplot[fill=yellow!80!black] coordinates {(Plan Fixe, 9.1) (Double DQN, 0.0) (MP-Sum, 0.0)};
\addplot[fill=green!70!black] coordinates {(Plan Fixe, 40.9) (Double DQN, 40.1) (MP-Sum, 80.8)};
\addplot[fill=purple!75] coordinates {(Plan Fixe, 9.1) (Double DQN, 18.3) (MP-Sum, 0.0)};
\legend{Phase 0, Phase 1, Phase 2, Phase 3}
\end{axis}
\end{tikzpicture}
\caption{Distribution temporelle relative des phases vertes (Graine 7).}
\label{fig:stacked_distribution}
\end{figure}

Les données de terrain révèlent un mécanisme fondamentalement différent du thrashing~:
\begin{enumerate}
    \item \textbf{Réfutation formelle du thrashing~:} MP-Sum n'effectue que \textbf{91 commutations} sur 1 heure, soit nettement moins que le plan fixe (159) et le DDQN (210). L'algorithme ne souffre donc d'aucune hyper-instabilité de commutation.
    \item \textbf{Phénomène de Famine de Phase (\emph{Phase Starvation})~:} La Phase~2 monopolise à elle seule \textbf{$80{,}8\,\%$ du temps vert total} ($2910\,\mathrm{s}$ sur $3600\,\mathrm{s}$). À l'opposé, la Phase~3 (tourne-à-gauche secondaire) subit une privation absolue de signal vert~: \textbf{$0{,}0\,\mathrm{s}$ de vert en 1 heure} !
    \item \textbf{Spillback et perte de capacité du réseau~:} Du fait de la géométrie asymétrique, les véhicules engagés dans les mouvements de la Phase~3 restent définitivement bloqués sur leur voie. La file s'allonge jusqu'à saturer la totalité de l'approche amont. Cette congestion physique empêche l'insertion de nouveaux véhicules générés par le simulateur, conduisant à un refoulement massif de \textbf{207 véhicules bloqués hors réseau} (Figure~\ref{fig:topology_graph}). C'est ce refoulement qui explique la chute de débit ($580$ vs $889$) et l'explosion mathématique du \emph{timeLoss} de MP-Sum.
\end{enumerate}

\subsection{Pourquoi l'Agent Double DQN Élimine la Famine}
Pourquoi la politique apprise par l'agent Double DQN évite-t-elle le piège de la famine dans lequel sombre Max-Pressure ?

La réponse réside dans la nature même des fonctions d'objectif~:
\begin{itemize}
    \item \textbf{Max-Pressure est une règle myope à un pas~:} Elle calcule $\arg\max_p \sum x_l(t)$ à l'instant présent. Dès lors que l'axe Est-Ouest présente structurellement un débit direct supérieur, la somme des files de la Phase~2 dépasse continuellement la file isolée de la Phase~3, condamnant cette dernière à ne jamais franchir le seuil d'activation.
    \item \textbf{Le Double DQN optimise l'horizon futur cumulé~:} Par le truchement du facteur d'actualisation $\gamma = 0{,}95$ et de l'équation de Bellman~(\ref{eq:target_double_dqn}), l'agent anticipe que laisser s'accumuler des véhicules immobiles sur la Phase~3 génère une croissance quadratique du temps d'attente cumulé global $W(t)$, dégradant les récompenses futures $r_{t+k}$ selon~(\ref{eq:reward}).
\end{itemize}

Par conséquent, l'agent DDQN apprend de manière endogène une stratégie d'arbitrage optimale~: il alloue $41{,}5\,\%$ du temps à la Phase~0, $40{,}1\,\%$ à la Phase~2, et déclenche de façon intercalaire des décharges courtes mais régulières sur la Phase~3 ($18{,}3\,\%$ du temps total, Figure~\ref{fig:stacked_distribution}). Cette politique de purge préventive maintient la fluidité de tous les axes et garantit un taux de blocage amont strictement nul ($N_{\text{block}} = 0$).

\section{Faisabilité de Déploiement et Conséquences Pratiques}
\label{sec:deployment}

Sur le plan de l'implémentation physique dans des environnements urbains contraints tels que Thiès, les résultats obtenus ouvrent des perspectives opérationnelles immédiates~:

\begin{itemize}
    \item \textbf{Frugalité d'instrumentation~:} L'état d'entrée de notre modèle ($\mathbf{s}_t \in \mathbb{R}^{13}$) ne requiert qu'une estimation macroscopique des longueurs de files et densités d'approches. De telles informations peuvent être extraites en temps réel par une caméra de vidéosurveillance grand angle couplée à un modèle de détection d'objets allégé (ex. YOLOv8n / Edge TPU) monté au sommet du mât de signalisation, pour un coût matériel inférieur à $500\,\$$, contrastant avec les dizaines de milliers d'euros nécessaires aux boucles SCATS/SCOOT.
    \item \textbf{Sobriété calculatoire~:} Avec seulement deux couches cachées de 64 neurones, l'inférence d'une décision par le réseau DDQN s'effectue en moins de $0{,}2\,\mathrm{ms}$ sur un micro-ordinateur monocarte (Raspberry Pi 4 ou NVIDIA Jetson Nano), satisfaisant amplement la contrainte de période de décision fixée à $5\,\mathrm{s}$.
\end{itemize}

\section{Conclusion et Travaux Futurs}
\label{sec:conclusion}

Cet article a présenté un benchmark expérimental exhaustif confrontant un contrôleur adaptatif basé sur l'apprentissage par renforcement profond (\emph{Double DQN}) aux standards théoriques et industriels du transport (Plan Fixe de Webster, Régulation Actionnée, et Max-Pressure) sous le simulateur microscopique Eclipse SUMO, appliqué à l'intersection réelle de Thiès (Sénégal).

L'analyse consolidée sur cinq graines stochastiques démontre qu'un agent Double DQN compact réduit de moitié le temps perdu moyen ($-50{,}7\,\%$) et de $62{,}2\,\%$ le temps d'attente à l'arrêt, tout en préservant le débit d'évacuation maximal. En outre, notre diagnostic microscopique apporte un éclairage inédit sur les limites pratiques de Max-Pressure sur carrefour asymétrique isolé~: l'effondrement de la règle gloutonne résulte d'une famine monopolistique de phase induisant un refoulement physique amont (\emph{spillback}), pathologie que l'agent Double DQN surmonte grâce à son optimisation à long terme de l'attente cumulée.

Les travaux futurs exploreront l'extension de ce formalisme à un réseau maillé multi-carrefours via des architectures d'attention sur graphes de type CoLight~\cite{wei2019colight} ou MAPPO, ainsi que l'évaluation de la résilience de l'agent en présence d'occlusions visuelles partielles ou de bruit de détection de caméras réelles.

\begin{thebibliography}{99}

\bibitem{haydari2022deep}
A.~Haydari and Y.~Y{\i}lmaz, ``Deep reinforcement learning for intelligent transportation systems: A survey,'' \emph{IEEE Transactions on Intelligent Transportation Systems}, vol.~23, no.~1, pp. 11--30, 2022.

\bibitem{webster1958traffic}
F.~V. Webster, ``Traffic signal settings,'' \emph{Road Research Technical Paper}, no.~39, Road Research Laboratory, HMSO, London, 1958.

\bibitem{lowrie1990scats}
P.~R. Lowrie, ``SCATS, Sydney Co-ordinated Adaptive Traffic System: A traffic responsive method of controlling urban traffic,'' \emph{Roads and Traffic Authority of New South Wales}, 1990.

\bibitem{hunt1981scoot}
P.~B. Hunt, D.~I. Robertson, R.~D. Bretherton, and R.~I. Winton, ``SCOOT: a traffic responsive method of coordinating traffic signals,'' \emph{TRRL Laboratory Report}, no.~1014, Transport and Road Research Laboratory, Crowthorne, 1981.

\bibitem{wei2018intellilight}
H.~Wei, G.~Zheng, H.~Yao, and Z.~Li, ``IntelliLight: A reinforcement learning approach for intelligent traffic light control,'' in \emph{Proc. 24th ACM SIGKDD Int. Conf. Knowl. Discovery Data Mining (KDD)}, 2018, pp. 2496--2505.

\bibitem{wei2019presslight}
H.~Wei, C.~Chen, G.~Zheng, K.~Wu, V.~Gayah, K.~Xu, and Z.~Li, ``PressLight: Decentralized reinforcement learning via localized reward for traffic signal control,'' in \emph{Proc. 25th ACM SIGKDD Int. Conf. Knowl. Discovery Data Mining (KDD)}, 2019, pp. 1916--1925.

\bibitem{chen2020toward}
C.~Chen, H.~Wei, N.~Xu, G.~Zheng, M.~Yang, Y.~Xiong, K.~Xu, and Z.~Li, ``Toward a thousand lights: Decentralized deep reinforcement learning for large-scale traffic signal control,'' in \emph{Proc. 34th AAAI Conf. Artif. Intell. (AAAI)}, vol.~34, no.~01, 2020, pp. 3414--3421.

\bibitem{zhang2020pdlight}
C.~Zhang, X.~He, and G.~Wang, ``PDLight: A deep reinforcement learning traffic light control algorithm with pressure and dynamic light duration,'' \emph{arXiv preprint arXiv:2009.13711}, 2020.

\bibitem{wei2021survey}
H.~Wei, G.~Zheng, V.~Gayah, and Z.~Li, ``Recent advances in reinforcement learning for traffic signal control: A survey of models and evaluation,'' \emph{ACM SIGKDD Explorations Newsletter}, vol.~23, no.~2, pp. 12--25, 2021.

\bibitem{zhao2024survey}
K.~Zhao, X.~Zhai, and D.~Guan, ``Deep reinforcement learning for traffic signal control: A comprehensive review of algorithms, environments, and benchmark metrics,'' \emph{IEEE Transactions on Intelligent Vehicles}, vol.~9, no.~3, pp. 4110--4128, 2024.

\bibitem{day2016detector}
C.~M. Day, H.~Li, L.~Richardson, J.~Howard, T.~Platte, J.~R. Sturdevant, and D.~M. Bullock, ``Detector-free optimization of traffic signal offsets with connected vehicle data,'' \emph{Transportation Research Record}, vol. 2620, no.~1, pp. 52--60, 2016.

\bibitem{varaiya2013max}
P.~Varaiya, ``Max pressure control of a network of signalized intersections,'' \emph{Transportation Research Part C: Emerging Technologies}, vol.~36, pp. 177--195, 2013.

\bibitem{lioris2016max}
J.~Lioris, R.~Pedarsani, F.~Y. Tascikaraoglu, and P.~Varaiya, ``Platooning for coupled car-following and maximum throughput at signalized intersections,'' \emph{Transportation Research Part C: Emerging Technologies}, vol.~71, pp. 295--312, 2016.

\bibitem{mnih2015human}
V.~Mnih \emph{et al.}, ``Human-level control through deep reinforcement learning,'' \emph{Nature}, vol. 518, no. 7540, pp. 529--533, 2015.

\bibitem{genders2016using}
W.~Genders and S.~Razavi, ``Using deep reinforcement learning for traffic light control,'' \emph{Transportation Research Record}, vol. 2620, no.~1, pp. 52--60, 2016.

\bibitem{van2016deep}
H.~van Hasselt, A.~Guez, and D.~Silver, ``Deep reinforcement learning with double Q-learning,'' in \emph{Proc. 30th AAAI Conf. Artif. Intell. (AAAI)}, 2016, pp. 2094--2100.

\bibitem{wei2019colight}
G.~Zheng, Y.~Xiong, X.~Zang, J.~Feng, H.~Wei, H.~Zhang, Y.~Li, K.~Xu, and Z.~Li, ``Learning phase coordination for large-scale traffic signal control,'' in \emph{Proc. 25th ACM SIGKDD Int. Conf. Knowl. Discovery Data Mining (KDD)}, 2019, pp. 1906--1915.

\end{thebibliography}

\end{document}
